\documentclass[12pt]{article}
\usepackage{graphicx} % Required for inserting images
\usepackage{enumitem}
\usepackage[a4paper, margin=2cm]{geometry}
\usepackage{amssymb}
\usepackage{amsmath}
\usepackage{bm}
\usepackage{mathptmx}
\usepackage{setspace}
\usepackage{stmaryrd}
\doublespacing
%\renewcommand{\footnotesize}{\small}

% --- APA CITATION SETUP ---
\usepackage{apacite}
\bibliographystyle{apacite}

\title{Essentialism and the Essence of \textit{to be Necessary}}
\author{Shin JeongBean}
\date{}
\begin{document}
\setlength{\parskip}{6pt}
\setlength{\parindent}{1em} 
\maketitle
%\vspace{-1cm}
% Your content starts here
\citeA{Fine1994-FINEAM-2} argued against the view that essence should be understood in terms of \textit{de re} modality. He suggested instead to understand metaphysical necessity by the essence of things. The idea to explain, especially to ground all metaphysical necessity in essence was followed by various philosophers albeit with some differences in details. I term all such accounts \textit{essentialism} on necessity.
\par
But if essence is not to be understood by modal notions, one might ask what makes essence so intimately related to necessity. In particular, some philosophers argued that essentialists cannot account for the reason why \textit{to be essential} is such that all facts obtaining essentially are necessary. If essence is partially grounded on necessity, essentialism infringes the asymmetry of ground; but if essence is non-modal, then there is a conceptual gap between what it is to obtain essentially, and what it is to be necessary. 
\par
In this paper I argue that there is no conceptual gap between non-modal conception of essence and necessity, if one takes metaphysical necessity to be the necessity which is essentially essentialist. In particular, the extensional equivalence between facts which obtain essentially and are metaphysically necessary can be taken as obtaining essentially of \textit{to be necessary}. I also argue that this move is forced upon essentialists. After some clarification on terminology in \S 1, \S 2 will be devoted to elaborate the challenge of gap, particularly raised against essentialists. In \S 3, I argue that the challenge should be answered by appealing to the essence of \textit{to be necessary}. In \S 4, I will clarify that my suggestion is legitimate, especially that it is not viciously circular. 
%------------------------------
\section{Introduction}
%------------------------------
By ``necessity'' I mean \textit{metaphysical} necessity if not otherwise stated. It is further supposed that all logically necessary or conceptually necessary facts are also metaphysically necessary. Similarly, it may be helpful to note that if there are varieties of grounding, by ``grounding'' I mean \textit{metaphysical} grounding. By ``essence'', I mean \textit{mediate and consequential} essence, following Fine (1995a). Suppose we have a clear idea of a strict and direct essence, which is immediate and constitutive. Then, the mediate essence of $F$ includes both the nature of itself and the nature of things that are involved in the nature of $F$; consequential essence of $F$ includes what logically follows from the nature of $F$. 
\par
Also crucial is the distinction between objectual essence and generic essence. On the one side there is objectual essence, the essence of a particular object, and the other side the essence of \textit{to be} such and such, that is, generic essence. Fine (2015) argued that the essence of \textit{to be} $F$ is different from the essence of the concept $F$. Suppose the objectual essence of the concept \textit{concept} includes the fact that \textit{concept} is a concept. Then it is a consequential essence of it that $\exists x (x \text{ is a concept})$. However, it seems quite implausible to think that it is by nature of what it is to be a concept that there is something satisfying it. That is, while $\exists x (x$ is a concept) is plausibly essential to the concept \textit{concept}, it would not be essential of \textit{to be a concept}. This, we may say, illustrates the difference between the objectual essence of concept and generic essence of \textit{to be a concept}. It would be clear that the essence of objects or things would be the objectual essence of that object. 
\par
I take the primary subjects of being essential, as well as being necessary to be facts (which not necessarily actually obtain), and take grounding to be the higher-order relation between facts. Yet it will not be very difficult for those who favor propositions, or first-order relational view of grounding translate the words used in this paper to their favored account. Sometimes I will use ``an essence of $F$'' to denote a fact essential to $F$, whereas ``nature'' is used strictly for the non-further-entity sense of essence (Lowe 2008). 
\par
I do \textit{not} use the two terms ``be essential of/to'' and ``obtain in virtue of the nature of'' (equivalent with ``obtain essentially of'') interchangeably, since I operate with non-factive notion of essence, where facts essential to whatever may not obtain. If a fact obtains essentially of something, it is essential to the thing in question; but not vice versa. Following Fine (2015), by $p \bm{\leftarrow}_{x_1, \dots ,x_n} F(x_1, \dots ,x_n)$ and $p \bm{\leftarrow} q$ I mean $p$ is a generic (mediate and consequential) essence of $G$ and $p$ is a generic essence of $q$ respectively. Objectual essence will be notated by using a \textit{rigid predicate} $\lambda x. (x = a \lor x = b \lor \dots)$, where $a, b, \dots$ are names for objects. To say $p$ is an objectual essence of $a$ I write ``$p \bm{\leftarrow}_x (\lambda y. y=a)x$''. For the notion of obtaining essentially of something, I borrow from the traditional notation \cite{Fine1995-FINTLO} and write ``$E_F p$'' for $p$ obtains in virtue of the nature of $F$, where ``$F$'' is possibly 0-placed. By convention, I sometimes write ``$E_a p$'' as an abbreviation of ``$E_{\lambda x. x=a} p$''. I also endorse the notation for generic grounding from Fine (2015), where $F(x_1, \dots, x_n) \bm{\leftarrow}_{x_1, \dots ,x_n} G(x_1, \dots ,x_n)$ means $F$ grounds $G$. Here the grounding is non-factive. 
\par  
By ``essence fact'' I mean the fact that whatever fact is essential of whatever. Furthermore, I assume that if whatever is essential of $F$, such an essence fact obtains essentially of $F$. That is, if $p \bm{\leftarrow}_{x_1,\dots,x_n} F(x_1,\dots,x_n)$, then $E_F (p \bm{\leftarrow}_{x_1,\dots,x_n} F(x_1,\dots,x_n))$. By ``realized essence fact'' I mean the fact that whatever fact obtains essentially of whatever, i.e., the fact that $E_F p$ for some $F$. I also assume that a realized essence fact obtains essentially: $E_Fp \supset E_F E_F p$ where ``$\supset$'' is used for material conditional to prevent ambiguity.
\par
By ``essentialism'', I denote the thesis that \textit{to be necessary} is grounded on \textit{to obtain essentially}, i.e., $E_F p \bm{\rightarrow}_p \Box p$. This would yield that for any necessary fact $p$, $\Box p$ is fully grounded by the fact that $p$ obtains essentially of whatever. This can be regarded as a proposed answer of accounting for the ultimate metaphysical source of necessity \cite{Hale2002-HALTSO-27}. The source of necessity is supposed to provide the metaphysical explanation of why a necessary fact is necessary. It does not seem likely that it must always also explain why the fact obtains. Sometimes that $p$ is necessary may be partially explained by that $p$ obtains. It is customary to explain the necessity of the identity between Hesperus and Phosphorus by the actual identity of them, and the necessity of identity\footnote{\citeA{Kripke1980-KRINAN-12}. Although it might be disputable whether this is a \textit{metaphysical} explanation, at least presumably there is no quick argument against it.}. In this case, the particular source for the necessity of $p$ would include $p$, which would be inadequate to explain $p$ itself. However, this cannot be the \textit{ultimate} source for necessity, for in explaining the necessity of the fact that Hesperus is Phosphorus we appealed to the necessity of identity. We might further ask why, then, every identity between an object and itself is necessary; and if it is again explained via appealing to the necessity of another fact $q$, we might ask why $q$ is necessary, and so on. 
\par
The chain of explanation might continue infinitely, or end at the point where we can give no further explanation for the necessity in question, i.e., where the necessity is primitive. Or, the chain will end if one can give the explanation for the necessity in question with non-modal source. The final non-modal destination would be, then, the \textit{ultimate} source of the particular necessity which was initially at stake. 
\par
Essentialism as I understand it, then, consists of two theses: that the ultimate source of \textit{all} necessities are some realized essence facts, and the source of necessity is the ground of the necessity. The first thesis is particularly strong. Some philosophers favor a weaker view: that some necessities have their ultimate sources on essence facts, but some other necessities do not\footnote{For example \citeauthor{Hale2013-HALNBA} (2013, Ch.6) holds that the necessity of objectual essence facts are primitive, and \citeA{Godman2020-MALEPA-3}, \citeA{Correia2021-COREMA-3} doubt that the necessity of logical truths has essentialist source.}. In this paper, I focus on the more ambitious essentialism which states all necessities are grounded in some realized essence facts. 
\par
The grounding thesis is explicitly endorsed by some philosophers\footnote{For example \citeA{Lowe2012-LOWWIT-2}, Hale (2013, Ch.6), and \citeA{Wallner2020-WALEEA-6}}. \citeA{correia2024non} suggests that the grounding thesis is not forced upon theorists who wish to seek the ultimate source of necessity in essence. He suggests that instead of grounding, one might hold that necessity is reduced to, or defined in terms of essence (as in \citeA{Fine1994-FINEAM-2} and \citeA{Fine1995-FINTLO}), or there is a generic identity between essence facts and necessities (as in \citeA{Correia2021-COREMA-3}). I admit that these are real possibilities, and by ``essentialism'' I \textit{exclude} accounts that do not involve the grounding thesis \footnote{For instance the account which holds that necessary facts and essential facts are extensionally equivalent, yet necessity is ungrounded \cite{Wilsch2017-WILSMP-3}. Similarly, \citeA{Correia2021-COREMA-3} hold that \textit{to be necessary} is generically identical with \textit{to be a logical consequence of a true generalized identity} (or its extensional correlate), thus the relation between them is symmetric, while grounding is thought to be asymmetric.} But it should be noted that there indeed are cases where the grounding thesis is implicit. For example, Kit Fine, when he defines \textit{to be necessary} in terms of \textit{to be essential of all objects}, he would have real constitutive definition in mind, which gives both grounding and essence claim \cite{Fine2015-FINUFF}. 
\par
Finally, some comments should be given for the restriction on \textit{objectual} essence which can be found occasionally. This qualification is explicit in some writings, for example \citeA{Fine1994-FINEAM-2}: ``The metaphysically necessary truths can then be identified with the propositions which are true in virtue of the nature of all \textit{objects} whatever'' (p.9, italics added); \citeA{Hale2013-HALNBA}: ``Metaphysical  necessities might plausibly be identified as those which hold in virtue of the nature of \textit{things} in general.'' (p.150, italics added). The general picture then would be that for any necessary fact $p$, we can find a rigid predicate $\rho$ such that $E_\rho p$ which grounds $\Box p$.
\footnote{
    Strictly speaking, in the logic proposed by \citeA{Fine1995-FINTLO}, for any predicate $F$ not necessarily rigid, given $E_F p$ we can find a rigid predicate $\rho$ such that $E_\rho p$. This is due to the axiom (II) (v): $F \subseteq G \supset (E_F A \supset E_G A)$. Intuitively, it says that if the intension of $F$ is a subset of the intension of $G$, any essence of $F$ is also an essence of $G$. Thus, given the rigid predicate of being one of all objects $\bigvee$, if $E_F p$, we could get $E_{\bigvee} p$. In this case, it seems to say for any necessary fact $p$ there is an essence fact such that $E_F p$ is at least extensionally equivalent to the original thesis proposed by Fine. However, the same axiom yields the metaphysically undesirable result that for any two intensionally equivalent predicates $F$ and $G$ $E_F p \Leftrightarrow E_G p$, regardless of their hyperintensional difference. Thus, I do not endorse the logic of \citeA{Fine1995-FINTLO} throughout the discussion. 
}
\par
In paradigms where essence of to be whatever is treated as a constitutive necessary condition (Fine, 2015), or as something generically identical with it \cite{Correia2017-CORGEA}, essences are straightforwardly not always necessary. For example, let ``$sg$'' stand for the fact that \textit{snow is green}. Then it is extremely plausible that $sg \bm{\leftarrow} sg$; yet certainly it is not necessary that snow is green. After all, it does not even actually obtain. This might seem to provide a motivation for the restriction on objectual essence, since we would not like to grant $E_{sg} sg$ and thus license $\Box sg$. However, $sg \bm{\leftarrow} sg$ simply does not grant $E_{sg} sg$ in our sense. As I mentioned earlier, the \textit{essence fact} about $F$ -- not the essence of $F$ -- obtains essentially of $F$. Thus, $sg \bm{\leftarrow} sg$ licenses only $E_{sg} (sg \bm{\leftarrow} sg)$; and as Fine (2015, p.299) suggested, $p \bm{\leftarrow} q$ yields $q \supset p$, thus it would mediately and consequentially yield $E_{sg} (sg \supset sg)$\footnote{
    In general, $p \bm{\leftarrow}_{x_1, \dots , x_n} F(x_1, \dots, x_n)$ would yield $\forall x_1 \dots \forall x_n (F(x_1, \dots, x_n) \supset p)$. 
}. This is acceptable, since we would have no problem accepting $\Box (sg \supset sg)$: necessarily, snow is green if snow is green. I would like to take one step further and protest that $E_{sg} (sg \bm{\leftarrow} sg)$ is \textit{grounded} on $sg \bm{\leftarrow} sg$, yet this is optional and perhaps contentious. 
\par
Thus, I see no particular reason for essentialists to think only objectual realized essence facts can ground necessity. It is even unclear how an essentialist would account for the necessity such as \textit{all bachelors are unmarried} without appealing to the generic essence of \textit{to be a bachelor}. By contrast, let ``$U$'' be the predicate ``\textit{to be unmarried}'' and ``$B$'' ``\textit{to be a bachelor}'', with generic essence we can have $Ux \bm{\leftarrow}_x Bx$, which yields $E_B \forall x (Bx \supset Ux)$, which would ground $\Box \forall x (Bx \supset Ux)$, i.e., it is necessary that all bachelors are unmarried. With such consideration, I will discard the restriction on \textit{objectual} essence throughout the discussion. Still, it is worth noting that all objectual essences \textit{regarding itself} obtain essentially of the object in question, in the mediate and consequential sense. For example, suppose $Fx \bm{\leftarrow}_x (\lambda y. y=a)x$. This yields $E_{\lambda y. y=a} \forall x ((\lambda y. y=a)x \supset Fx)$, and by $\beta$-reduction $E_{\lambda y. y=a} \forall x (x=a \supset Fx)$, and by the rule of identity $E_{\lambda y. y=a} Fa$. 

%------------------------------------
\section{The Challenge of Gap}
%--------------------------------------

Obviously, all essentialist accounts should not take \textit{to obtain essentially} to be understood in terms of \textit{to be necessary}. For example, suppose for $Fa$ to obtain essentially of $a$ just is for $a$ to necessarily have $F$. Then to say the necessity of $Fa$ is explained by $Fa$ obtaining essentially just amounts to saying that $Fa$ is necessary because $a$ necessary has $F$. In this case essentialism results in blatant circularity. 
\par
On the other hand, it is the thesis of essentialism that for any $F$ and $p$, $E_F p$ grounds $\Box p$. But then one might ask if essence should not be understood in terms of modality, i.e., has a non-modal nature, how one can ensure such a close relationship between \textit{to obtain essentially} and \textit{to be necessary}. The challenge is often articulated in terms of extensional relation between them. The grounding thesis implies what \citeA{Mackie2020-MACMBE} calls ``Necessity Principle (NP)'': 
\begin{itemize}
    \item[(NP)] $\forall p (E_F p \supset \Box p)$\footnote{where ``$F$'' is any $n$-place predicate for any $n \in \mathbb{N}_0$}
\end{itemize}
and it is challenged that if essence is non-modal, essentialists cannot account for why (NP) holds. Of course, essentialists may contend that (NP) obviously holds. However, the problem I want to discuss here is not epistemic but metaphysical.\footnote{\citeA{Correia2021-COREMA-3} gave an excellent argument in reply to the epistemic challenge, and it seems that the premise of the argument can be weakened so that any essentialists who hold essence is intimately related to identity (not necessarily defined in terms of identity) can utilize.} Even if essentialism is correct, (NP) would not be a fundamental fact of reality\footnote{I assume that only ungrounded facts can be fundamental, that is, the `minimalist' worldview of \citeA{Fine2001-FINTQO}.}. There seem at least two reasons for essentialists to believe the nonfundamentality of (NP). First, since (NP) involves universal quantification, it is plausibly grounded on each instance of $p$ satisfying $E_F p \supset \Box p$, which is plausibly further grounded on either $\lnot E_F p$ or $\Box p$. Second, essentialism implies that \textit{to be necessary} is grounded, and therefore nonfundamental; whereas if (NP) is fundamental, assuming the principle of purity any constituent of (NP) including \textit{to be necessary} (viz., `$\Box$') would also be fundamental, thus contradiction.  
\par
Thus, it would make sense to ask for the metaphysical explanation for (NP). But then it is unclear how non-modal conception of essence can give metaphysical explanation of why (NP) holds. (NP) would have a solid metaphysical basis if \textit{to obtain essentially} is simply defined in terms of necessity. Yet as we saw before, this is not something essentialists can admit. The predicament can be illustrated as follows. To avoid circularity, essentialists should not take essence, or more specifically \textit{to obtain essentially} to have modal nature. But to provide metaphysical basis for (NP), it seems that they should think there is a modal import in essence. But it is impossible to say essence is both non-modal and modal. As Mackie writes: how can essentialists ``deliver a modal rabbit out of a non-modal hat''? (Mackie, 2020, p.252, also see \citeauthor{Leech2018-LEEEAM}, 2018 and Romero, 2019)
\par
Preexisting attempts by essentialists to provide a metaphysical basis for (NP) faced strong objections. One common understanding of essence among essentialists is by \textit{real definition}: roughly, it is the thesis that the constitutive and immediate essence of certain thing or \textit{to be such and such} is captured by their real definition\footnote{See, for example Fine (1994), \citeA{Lowe2012-LOWWIT-2}, or Hale (2013)}. Hence, if some $x$ lacks the essential property of an object $\alpha$, it means that it fails to satisfy the definition of $\alpha$, thus we can conclude that $x \neq \alpha$. Hale tried to argue that this conception of essence already grants (NP). Suppose $\Phi \alpha$ is an objectual essence of $\alpha$. Now suppose it is possible that $\lnot \Phi \alpha$. It implies that there is a $\beta$ such that $\alpha = \beta \land \lnot \Box \Phi \beta$. But $\beta$ cannot be identical to $\alpha$ in the world such that $\lnot \Phi \beta$. For, it does not satisfy the condition to be identical to $\alpha$, which is $\Phi$. Therefore by the necessity of identity $\lnot \alpha = \beta$, contradiction. Since contradictions are impossible, we can conclude that it is not possible that $\lnot \Phi \alpha$. In other words, if $E_{\lambda x. x = \alpha} \Phi \alpha$, then $\Box \Phi \alpha$ (Hale 2013, p.133). 
\par
But it has been pointed out that definitional conception of essence cannot guarantee (NP) at least in the way Hale proposed (\citeauthor{Romero2019-ROMMIN} 2019, pp.125--7; \citeauthor{Leech2021-LEEFET-2} 2021, pp.893-5). For the proposed argument to succeed, it must be assumed that $\Phi$ is \textit{necessarily} the definition of $\alpha$. That is, either $\Box E_{\alpha} \Phi \alpha$ or $\Box \Phi \alpha$ is required. But what is assumed is only that $\Phi$ is \textit{actually} the condition of the identity. And to hold that the assumption $E_{\alpha} \Phi \alpha$ implies $\Box \Phi \alpha$ simply begs the question.\footnote{Of course, Hale himself would wish to reply that the necessity of $E_{\alpha} \Phi \alpha$ is not to be explained by further essence. Rather, it is a fundamental necessity that does not permit further explanation (Hale 2013, p.158). However, this move, whether it is convincing or not, is not permitted to essentialists in our sense.}
\par 
Alternatively, \citeA{Correia2017-CORGEA} suggested that essence should be understood in terms of generalized identity. Let ``$\phi \equiv \psi$'' mean that $\phi$ is generically identical to $\psi$. Then $q$ is a full essence of $p$ iff $p \equiv q$, and $r$ is a partial essence of $p$ iff there is a $q$ such that $p \equiv (r \land q)$ (pp.649-50). And to understand essence in terms of generalized identity may seem to offer better prospects for advocating essentialism\footnote{As far as I know, there are no actual essentialists who argue for essentialism in our narrow sense with such a conception of essence. The allegedly failed argument below is rather constructed by challengers of essentialism (Romero, 2019; Leech, 2021).}, since ``[a]s with objectual identity, generalized identity is tightly linked to metaphysical necessity'' (\citeauthor{Correia2017-CORGEA} 2017, p.656). Here $p \equiv p$ is a theorem\footnote{Since they take generalized identity to be reflexive (p.645)} and therefore in any normal modal logic it follows that $\Box (p \equiv p)$. And given $p \equiv (q \land r)$, which means that $q$ is a partial essence of $p$, by substitutability we get $\Box (p \equiv (q \land r))$. That is, assuming $q$ is a partial essence of $p$, it follows that $q$ is a partial essence of $p$ \textit{necessarily}.
\par
This actually proves something stronger than (NP): namely, $E_F p \supset \Box E_F p$. However, since whatever obtains essentially of something also obtains \textit{simpliciter}, this consequence implies (NP). The problem here is that the proof presupposes $\Box (p \equiv p)$: that the generalized identities are necessary. Then what could be the source of the necessity of $(p \equiv p)$? It has been argued that the necessity of $\Box (p \equiv p)$ cannot be explained by the essence of \textit{to be generically identical} (Leech 2021, pp.901-2; also see Romero 2019, p.138). For example, suppose that the necessity is explained by that \textit{for two things to be generically identical is partially to be in a reflexive relation}. Grant, but what explains the necessity of the partial generalized identity between generalized identity and reflexive relation? Here we cannot appeal to the necessity of generalized identity, for that exactly is the supposed explanandum. Therefore, to appeal to the essence of generalized identity does not solve the problem.
\par
Leech seems to conclude that in explaining the necessity of the generalized identity ``\textit{one should not appeal to the essence of anything to account for this necessity}''. She further considers this constraint as ``the crucial and general objection'' on the necessity of facts obtaining essentially (p.902). 
\par 
I disagree. I think that Leech overlooks a remaining possibility: that all generalized identities are necessary because of the essence of \textit{to be necessary}. Furthermore, in my view, this also applies directly to (NP). I agree that any non-modal specification of essence cannot be used to metaphysically guarantee (NP). However, every fact obtaining essentially is necessary, not because of the nature of what it is to obtain essentially but because of the nature of what it is to be necessary. In the next section, I first argue that essentialists must take (NP) to obtain essentially of \textit{to be necessary}, and show how adopting this thesis helps essentialists respond to the challenge of gap. 

%---------------------------------
\section{The Essentialist Essence of \textit{to be Necessary}} %----------------------------------

Essentialists who take (NP) to hold would also wish to take (NP) to hold \textit{necessarily}. There are various possible motivations behind. If an essentialist takes side on \textit{grounding necessitism}, viz., the thesis that the grounds necessitate the grounded, it would be a natural consequence that (NP) is necessary. Due to the popularity of necessitism, this might already provide good reason for essentialists in general to adopt the necessity of (NP)\footnote{Furthermore, necessitism is extremely natural for essentialists who adopt a `unified foundation' of essence and ground.}. But even if one favors contingentism, viz., the thesis that the grounds do not necessitate the grounded, may still take (NP) to hold necessarily in order to avoid conflict with plausible metaphysical assumptions. To reason in terms of modal logic without giving metaphysical import: if (NP) is not necessary (in the actual world $@$) there is a possible world accessible from $@$ where $\exists p (E_F p \land \lnot \Box p)$. And under essentialism the witness of $p$ either does not obtain essentially in $@$ or it does and therefore $\Box p$ holds in $@$. In the latter case, we straightforwardly get a counterexample to the S4 principle: $\Box p \supset \Box \Box p$, which is commonly thought to hold for metaphysical modality. In the former case, since essentialists also hold that $E_F p \supset \Box E_F p$\footnote{Immediate from $E_F p \supset E_F E_F p$ and (NP).}, insisting S4 it follows that all realized essence facts obtaining in $@$ obtain in all possible worlds; but in some non-actual possible worlds there obtain additional realized essence facts. But there does not seem to be any good reason to think that the actual world has the fewest realized essence facts among all possible worlds, if there can be a discrepancy of realized essence facts between worlds at all.  
\par 
Thus, in any case, essentialists would admit that (NP) is necessary. And given essentialism, to say that (NP) is necessary amounts to saying that (NP) obtains essentially of whatever. But of what would it obtain essentially, if it does at all? 
\par
It might be more promising to search for the subject of which (NP) obtains essentially of by inspecting (NP) itself. Notice the formalization of (NP) is articulated by a universal quantification over facts, an operator $E$ which reads ``to obtain essentially of'', a material conditional, and a necessity operator $\Box$ which reads ``to be necessary''. But it would be obvious that (NP) would have nothing to do with the nature of the material conditional, so we are left with three candidates.\footnote{One might insist that I should also consider a polymorphic existential quantifier over $n$-place predicates for $n \in \mathbb{N}_0$. But it is contentious whether such a quantifier is desirable for the object-language; and even if it is desirable, I do not think it makes sense to think (NP) obtains essentially of whatever that object-language term is associated with. At least I fail to make sense of such an idea.} It might seem equally farfetched to think that (NP) obtains essentially of the logical object which is associated with the universal quantifier over facts. It might sound better to say (NP) obtains essentially of \textit{all facts}. The general idea might be that, for each individual fact $p$ it obtains essentially of $p$ that $E_F p \supset \Box p$; and the totality of facts, if somehow given that the totality of facts contains \textit{all} facts, it would obtain essentially of the totality of facts that for every fact, if it obtains essentially then it is necessary, i.e., (NP). 
\par
But could (NP) obtain essentially of yet something else? Well, trivially, if (NP) obtains essentially of one of three candidates mentioned in the previous paragraph, it will also obtain essentially mediately of anything whose essence involves it. But let us exclude such cases from consideration due to their metaphysical insignificancy. Still, it seems one possible line of thought is that, although (NP) is not fundamental itself, it might obtain essentially of fundamental metaphysical laws, or something similar.
\par
In sum, we are left with four suggestions: either (NP) obtains essentially of \textit{to obtain essentially}, of all facts, or of some fundamental metaphysical laws, or obtains essentially of \textit{to be necessary}. I argue in detail that every option except the last fails. Then I proceed to show how the final option helps to answer the challenge of gap on behalf of essentialists. 
%------------------------
\subsection{Taking (NP) to obtain essentially of \textit{to obtain essentially}}


It is worth noting that it is already suggested by \citeA{Wilsch2017-WILSMP-3} and \citeA{Wallner2020-WALEEA-6} that (NP) or something stronger than that obtains essentially of \textit{to obtain essentially}, although the former would not count as essentialism in our sense. Wilsch directly takes (NP) to be the essence of essence, whereas Wallner and Vaidya take both $E_F p \supset \Box E_F p$ and $E_F p$ \textit{can explain} $\Box p$ essential to the operator $E$. Since $E_F p \supset \Box E_F p$ entails (NP), it would be fair to say both takes (NP) as a mediate and consequential essence - which I simply say `essence' - of \textit{to obtain essentially}.\footnote{
Wallner and Vaidya explicitly endorse that ``modality is grounded in definitional essences.'' (Wallner and Vaidya 2020, p. 431). They also make it explicit that they intend to follow Hale in giving non-transmissive explanation for the source of necessity in terms of ground. However, in case of Wilsch, he explicitly states himself as a modal primitivist, not an essentialist. Although he takes essence to give the ``source'' of metaphysical necessity, here ``[t]he sources of necessity are the modal joints of nature.'' (Wilsch 2017, p.429), which determine the extension of various kinds of necessity and account for the naturalness of them. Essence, then, explains why metaphysical necessity is a significant kind of primitive necessity, why it ``carves the nature at its joints''. This is apparently a different enterprise from the essentialism being discussed in this paper.} Following Wilsch, let us term the thesis that (NP) obtains essentially of \textit{to obtain essentially} ``Modal Essence of Essence'', ``(MEE)'' for short. That is: 
\begin{itemize}
    \item[(MEE)] $E_E \forall p (E_F p \supset \Box p)$
\end{itemize}
\par
In response, \citeA{Romero2023-CAROEN} argued that (MEE) results in the claim that $p \land q$ grounds $p$ which infringes the irreflexivity\footnote{To be precise, he appeals to the irreflexivity of \textit{metaphysical explanation} instead of grounding, since Wilsch is not committed to the thesis that essence grounds necessity. I focus on grounding to keep the discussion compact.}. His argument is convincing if one adopts the conception of essence in terms of generalized identity. Suppose that $E_F p \equiv_p (\Box p \land q)$: to obtain essentially just is to be necessary and $q$. This would yield that $E_{E} \forall p(E_F p \supset (\Box p \land q))$ and consequentially $E_{E} \forall p(E_F p \supset \Box p)$ if we follow \citeA{Correia2017-CORGEA}. Now, since two generically identical statements would be substitutable even in some hyperintensional contexts (\citeauthor{Dorr2016-DORTBF-2}, 2016; also see \citeauthor{Rayo2013-RAYTCO-4}, 2013), if $p$ grounds $q$, $r$ such that $p \equiv r$ would also ground $q$. A brief example: suppose $x$'s not being the subject of a marriage loan is grounded by $x$'s being a bachelor. Since to be a bachelor just is to be an unmarried adult male, $x$'s not being the subject of a marriage loan would also be grounded by $x$'s being an unmarried adult male. Similarly, given $E_Fp \equiv_p \Box p \land q$, the essentialist thesis that $E_Fp$ grounds $\Box p$ is equivalent to the claim that $\Box p \land q$ grounds $\Box p$. And that apparently does not seem satisfactory, therefore (MEE) is inconsistent with essentialism. 
\par
However, it is not obvious whether this argument works without \citeauthor{Correia2017-CORGEA}' view on essence. The argument crucially relies on the premise that if $q$ is the full essence of $p$, ``$q$'' must be able to substitute ``$p$'' in grounding context. This is not intuitively obvious for two reasons. First, it is not very clear there is a metaphysically significant notion associated with ``full essence''. Second, even we grant the notion of full essence, one can still contend that the full essence of $p$, say $q$, is intimately related to, but yet not (generically) identical to $p$. Then we would at least require a substantial theory of essence and ground to insist that if $p$ grounds $r$, $q$ also grounds $r$. 
\par
But there is a general argument for the incompatibility of essentialism and (MEE). Let us first introduce the notion of essential dependence: $x$ essentially depends on $y$, in Fine's term: ``if $y$ is a constituent of a proposition that [obtains] in virtue of the identity of $x$”. (\citeauthor{Fine1995-FINXD} 1995b, p.275; the term ``essential dependence'' follows \citeauthor{Correia2008-COROD} (2008)) Therefore from $E_E \forall p(E_F p \rightarrow \Box p)$, it follows that \textit{to obtain essentially} depends on the constituent $\Box$, that is, \textit{to be necessary}. However, since essentialists also hold that $E_F p \bm{\rightarrow}_p \Box p$, they also hold that \textit{to obtain essentially} does not depend on \textit{to be necessary}, contradiction. Of course essential dependency is not antisymmetric. However, what essentialists need is that \textit{to be} $C$ is strictly prior\footnote{
    Following \citeA{Fine1995-FINXD}, I use the term ``dependency'' for the fundamentality relation that does \textit{not} imply antisymmetry (p.282-3). If $X$ depends on $Y$, one only can infer that $X$ is \textit{not prior} than $Y$, for the possibility of $Y$ being dependent on $X$ is still open. The term ``prior to'' is reserved for the antisymmetric fundamentality relation.
}, since standard understandings of ground (cf. \citeauthor{Schaffer2009-SCHOWG} 2009, \citeauthor{Fine2012-FINGTG} 2012) take the ground strictly prior, or more fundamental to the grounded. Thus, essentialism and (MEE) are incompatible.
\par
The appeal to essential dependence may seem somewhat awkward, for I am operating with the consequential sense of essence, whereas the notion of essential dependence is usually formulated in terms of constitutive essence for good reasons. For example, assume the standard definition of biconditional `$\Leftrightarrow$' can be also read as a generic grounding claim: $((p \supset q) \land (q \supset p)) \bm{\rightarrow} (p \Leftrightarrow q)$. The obvious implication would be that the higher-order relation $\supset$ is strictly prior to the higher-order relation $\Leftrightarrow$. Yet since $p \Leftrightarrow p$ is a logical truth, it would be a consequential essence of anything, including $\supset$. And if consequential essence also yields essential dependence, it turns out that $\supset$ depends on $\Leftrightarrow$, contradiction.\footnote{
    \citeA{Fine1995-FINXD} tries to solve this problem by introducing the additional condition: ``$x$ depends upon $y$ just in case $y$ cannot be generalized out of the consequentialist essence of $x$'' (p.278), where $y$ can be `generalized out' of the essence of something if all occurrence of $y$ in the essence can be substituted by any object and still belongs to the essence. This perhaps can solve the problem regarding the dependence between \textit{objects}. Nevertheless, it does not seem to work regarding dependence outside first-order objects. What I proposed seems to be a counterexample: $\Leftrightarrow$ in the consequential essence of $\supset$ cannot be generalized out, for substituting into some two-place relation over propositions or facts would not result in logical 
    truths. Also see \citeA{VogtForthcoming-VOGUEH}
}
\par
Now, if (NP) is proposed as obtaining constitutively essential of \textit{to obtain essentially}, the aforementioned worry does not occur - and my argument still works fairly well. And even in some case where it is constitutive it may not affect the validity of the argument. For example, essentialists in fact must accept not only the conditional but also the biconditional $\forall p (E_F p \Leftrightarrow \Box p)$. Then given the motivation, one might want to say it is in fact the biconditional which obtains essentially of \textit{to obtain essentially}, and (NP) is merely consequential. I find this proposal reasonable. Nevertheless, my argument \textit{still} works fairly well, for ``$\Box$'' occurs in the biconditional as well.
\par
So we must examine \textit{how} (NP) would obtain consequentially essentially. The following, although not mutually exclusive, exhausts the logical possibility given (MEE): (1) (NP) follows from logic alone; (2) the set of facts obtaining essentially of \textit{to obtain essentially} is inconsistent; (3) $\forall p \lnot E_F(p)$ follows from what obtains constitutively essentially; (4) $\forall p \Box p$ follows from what obtains constitutively essentially; (5) constitutive essence includes $\Box$; (6) constitutive essence includes $\Diamond$.\footnote{I have classical S5 logic in mind; but it is also exhaustive even if one endorses S4 or intuitionist logic.} 
\par
(1) -- (4) are untenable, whereas (5) and (6) shows that \textit{to obtain essentially} still depends on \textit{to be necessary} or something similar. (1). It would be clear that (NP) does not follow from some customary logics - classical first-order logic, normal modal logic, and so on. If (NP) is a logical fact in that sense, the challenge of gap would not arise at all, for we do not generally ask for the metaphysical explanation for a logical truth even if the explanation can be given. Of course, we might construct a logic where (NP) follows logically. But then we will be confronted with the problem of philosophy of logic that why should we adopt such a logic, and in all likelihood the philosophical motivation would be essentialism itself. Then the opponent can raise the challenge of gap again, which will bring us back here once again. Therefore, one cannot appeal to (1) to escape the priority problem. 
\par
(2). Since what obtains in virtue of the nature of what would obtain \textit{simpliciter}, taking (2) means taking the world to be inconsistent. I do not think this is a price many philosophers want to pay, and those who are willing to pay the price would typically refuse to say we can infer everything from inconsistent premises. (3). Together with essentialism this means that there are no sources of necessity, and therefore no necessary facts. It is intuitively absurd, and I do not think the task of making sense of the proposal and refuting it would be adequate for the purpose of the paper. At least dialectically, I expect one who enters the debate of the source of necessity would not hold that there are no necessary facts. The same consideration applies to (4), which means that all facts are necessary.  
\par
(5). It would be obvious that if the constitutive essence of \textit{to obtain essentially} includes $\Box$ it follows that \textit{to obtain essentially} indeed depends upon \textit{to be necessary}. And as I argued, this undermines the initial essentialist account. To show (6) still results in undesirable consequence is not straightforward. One easy way out, is of course to take $\Diamond$ to depend on $\Box$, so that you cannot get (6) without (5). But perhaps one would prefer to take $\Box p$ to be defined as $\lnot \Diamond \lnot p$, or two operators to be on a par in regard to fundamentality. Throughout the paper, I implicitly assumed that $\Box$ is the fundamental modal operator, and essentialism is an account of the source of \textit{necessity}, not other kinds of metaphysical modality. However, this assumption was made only because it facilitates the discussion. It is certainly possible for the essentialist account to be formulated in terms of possibility: the source of \textit{possibility} of any possible fact $q$ is explained by $\lnot E_F(\lnot q)$. If $\Diamond$ is more fundamental than $\Box$, essentialism probably should undergo such a modification. But once the modification is made, (6) undermines the account, for it implies that \textit{to obtain essentially} is dependent upon \textit{to be possible}. 
\par
Hence, endorsing (MEE) would in any case make \textit{to obtain essentially} essentially dependent on \textit{to be necessary}. Still it might be argued that this does not preclude \textit{to obtain essentially} grounding \textit{to be necessary}. The preclusion is by the following reasoning. (i) If \textit{to be F} essentially depends on \textit{to be G}, \textit{to be F} is not prior to \textit{to be G}; (ii) If \textit{to be F} grounds \textit{to be G}, \textit{to be F} is prior to \textit{to be G}. Conclusion: if \textit{to be F} essentially depends on \textit{to be G}, \textit{to be F} cannot ground \textit{to be G}. The criticism is that for the inference to be valid, ``prior to'' in (i) and (ii) should express the same metaphysical relation, which does not seem very likely \cite{Vogt2026-VOGDAD-2}. 
\par
But to draw the conclusion one need not hold that (i) and (ii) operates with the identical notion of priority. To show that the priority cannot head towards opposite directions would suffice. This would be done, for example, if we show that both grounding and essence gives constitutive condition, where presumably constitution cannot be circular \cite{Fine2015-FINUFF}. Also, one can argue that the real definition - which is sometimes taken to be the constitutive and immediate essence - of whatever can be formulated in terms of grounding (\citeauthor{Rosen2015-ROSRD-2} 2015, pp.198-9\footnote{Notice that the final proposal in p.200 is strikingly similar with \citeA{Fine2015-FINUFF}: a real definition giving a necessary and sufficient condition must involve both essence and ground.}). I am inclined to think the account of Fine (2015) is largely correct, but to argue for the macroscopic metametaphysical theory would certainly exceed the scope of the paper. Still, some direct arguments for the opposite-heading of essence and ground should be dealt. 
\par
\citeA{Vogt2026-VOGDAD-2} suggests that, assuming Aristotelian account of property, there can be cases where $E_a Pa$ but $Pa$ grounds \textit{property $P$ exists} (p.519). Vogt seems to regard this being the case that $a$ being both prior (by grounding) and not prior (by essential dependence) to $P$. However, it should be noticed that `$P$' is used once as a name for a first-order object, once as a predicate. Thus, what the grounding statement and the essence statement jointly imply is actually that $a$ is not prior (by essential dependence) to \textit{to be} $P$; where $a$ is prior (by grounding) to the \textit{property} $P$. This is clearly not a case of two priority relations heading to the opposite direction. 
\par
Vogt (2026, in press) also cites various specific cases where one might contend that $Fx \bm{\rightarrow}_x Gx$ but $E_F \forall x(Fx \supset Gx)$ (see also \citeA{Correia2018-CORODG}). Of course, she is cautious enough not to declare any of the particular example to be definitive. However, it is suggested that the candidate examples, taken together, have enough force to make one take such possibility seriously. I disagree, for it seems that, even though that the grounding claim and the essentiality of the conditional is correct, those putative examples are subject to a uniform treatment. That is: the conditional should obtain essentially of the grounded, not the ground. For example, Vogt holds in both paper that \textit{a promises that p} grounds \textit{a has the obligation to p}, whereas it seems that it obtains essentially of \textit{to promise} that (one promises that $p$ $\supset$ one has the obligation to $p$). One of course may deny that there is such an intimate connection between \textit{to promise} and \textit{to have obligation}. If there is, my treatment is to deny the essence claim, and alternatively say it obtains essentially of \textit{to have obligation} that (one promises that $p$ $\supset$ one has the obligation to $p$).\footnote{This is exactly parallel to the goal of this section. Essentialist would agree that $E_F p$ grounds $\Box p$. But if one also have to think that the conditional $E_Fp \rightarrow \Box p$ should be essential, then it should be essential of the grounded, that is, \textit{to be necessary}, not of the ground, that is, \textit{to obtain essentially}.}\footnote{One might criticize such a treatment as \textit{ad hoc}. It perhaps is in some sense, but notice this does not increase the theoretical complexity of priority, but reduces it. If the priority by grounding and priority by essential dependence was so fundamentally different so that it permits opposite-heading, it seems that there should be two fundamentally different account of what metaphysical priority is, perhaps even two different kinds of fundamental reality, which should be altogether compatible. That not being present, I would regard my move justified.} 
\par
Therefore, insofar as essentialists try to ground \textit{to be necessary} - alternatively, \textit{to be possible} - in \textit{to obtain essentially}, they should not think (NP) obtains in virtue of the nature of \textit{to obtain essentially}. 

%------------------------
\subsection{Taking (NP) to obtain essentially of fundamental laws, or all facts}
%--------------------------

Extremely similar considerations would lead us to reject the suggestion that (NP) would obtain essentially of some fundamental laws, that is:
\begin{itemize}
    \item[(MEL)]: $E_L \forall p(E_F p \supset \Box p)$, where $L$ is the (totality of) law of which (NP) allegedly obtains essentially of.
\end{itemize}

Recall that I argued that (NP) is not fundamental in \S 2. Then we can utilize again the notion of essential dependence, and argue that (MEL)\footnote{Modal Essence of Laws.} implies that $L$ depends on (NP), and therefore (NP) is prior to those laws or is as fundamental as them. But since nothing can be prior to the laws since they are assumed to be fundamental,  (NP) would as fundamental as the fundamental laws, i.e., fundamental. This contradicts with the observation that (NP) would not be fundamental, especially if essentialism holds. 
\par
Again the argument in the previous paragraph is awkward since it operates with the consequential notion of essence. Still, the list (1)--(6) is still exhaustive even if we substitute the subject of obtaining essentially of into the fundamental laws in question, and exactly same argument can apply. To repeat: given (MEL), either (1)$'$ (NP) follows from logic alone, (2)$'$ the set of facts obtains essentially of $L$ is inconsistent; (3)$'$ $\forall p \lnot E_F(p)$ follows from what obtains constitutively essentially of $L$; (4)$'$ $\forall p \Box p$ follows from what obtains constitutively essentially of $L$; (5)$'$ constitutive essence includes $\Box$; (6)$'$ constitutive essence includes $\Diamond$. (1)$'$ is obviously not the case. (2)$'$ implies that the actual world is inconsistent; and if one can get along with this consequence, one will not think (2)$'$ would license (NP) to obtain essentially of $L$ automatically. (3)$'$ says nothing is necessary given essentialism, and (4)$'$ straightforwardly says everything is necessary. They are implausible \textit{per se}, and moreover anyone who enters the debate of the source of necessity would not wish to commit to either of these two theses. Finally, (5)$'$ and (6)$'$ yields the incompatibility between essentialism (or its suitably modified version) and (MEL).
\par
And it is again the consideration of dependency that undermines the idea that (NP) might obtain essentially of \textit{all} facts. The idea suggests that for each fact $p$, $E_F p \supset \Box p$ obtains essentially of $p$. Now suppose a witness $q$ which is a fundamental fact. By essential dependence, it follows that $\Box$ is fundamental; whereas by the grounding thesis of essentialism it follows that $\Box$ is not fundamental. 
\par
One might criticize that it is unfair to make the assumption that if (NP) obtains essentially of all facts, every particular fact should have necessity-related essence. The problem I raised will not occur if, only some nonfundamental facts have necessity-related essence, and these limited cases are sufficient to grant (NP) to obtain essentially of the totality of facts. I find such scenarios implausible. In those scenarios, not only selected nonfundamental facts should guarantee, by their nature, that for numerous facts, especially fundamental ones, if they obtain essentially then they are necessary; it should also guarantee that those facts are \textit{all the facts out there}. I simply fail to imagine any putative cases where the latter can be achieved. \footnote{It is indeed possible that (NP) is \textit{essential} to some nonfundamental facts. In fact, it is plausible that (NP) is essential to (NP) itself. But we cannot infer that (NP) \textit{obtain essentially} of whatever from it.}  
\par
But there is a final suggestion worth discussing. Instead of appealing to the essence of \textit{all facts}, we may appeal to the essence of \textit{to be a fact}. That is:
\begin{itemize}
    \item[(MEF)] $(E_F X \supset \Box X) \bm{\leftarrow}_X \textit{Fact}(X)$.
\end{itemize}
(MEF)\footnote{Modal Essence of Facts.} does not directly grant $E_{\textit{Fact}}$(NP), but rather $E_{\textit{Fact}} \forall X (\textit{Fact}(X) \supset (E_F X \supset \Box X))$. But since it is extremely plausible that it is either the case that $E_{\textit{Fact}}\forall p \ \textit{Fact}(p)$ or $E_{\forall_t} \forall_t p \ \textit{Fact}(p)$ where $\forall_t$ is the quantifier quantifying over facts, it can be supposed that under the suggestion either $E_{\textit{Fact}}$(NP) or $E_{\forall_t}$(NP) would hold. 
\par
However, (MEF) it is not metaphysically convincing. First, it is intuitively plausible that if the primary subject of being necessary or obtaining essentially are facts, \textit{to be a fact} would be prior to \textit{to be necessary} and \textit{to obtain essentially}. Among a variety of facts, only a proper subclass of them would obtain essentially (and assuming essentialism, therefore necessary). I submit that this gives at least \textit{prima facie} reason to think \textit{to be a fact} is prior to \textit{to obtain essentially}, especially if essentialists take a \textit{fact} being necessary is grounded on a \textit{fact} obtaining essentially. For an analogy, it is more plausible to think \textit{to be a natural number} is prior to \textit{to be prime}, especially if we take a \textit{natural number} being prime is grounded on a \textit{natural number} not being a product of two natural numbers less than itself. At least, we definitely would not want to take that if a natural number is not a product of two natural numbers less than itself then it is prime as an essence of \textit{to be a natural number}. 
\par
Second, ??? [I feel the first argument is somewhat weak, but I failed to come up with a convincing second argument]



%------------------------
\subsection{Taking (NP) to obtain essentially of \textit{to be necessary}}
%--------------------------
The only remaining possibility, then, is to say that (NP) obtains essentially of \textit{to be necessary}, that is:
\begin{itemize}
    \item[(EEN)] $E_\Box \forall p (E_F p \supset \Box p)$ 
\end{itemize} 
It would be clear that in this case there is no problem with regard to fundamentality, since it follows that \textit{to be necessary} is not prior to the constituent \textit{to obtain essentially}, which is harmonious with the grounding claim that $E_F p$ grounds $\Box p$. 
\par
But besides the metametaphysical harmony, we should also examine whether (EEN)\footnote{Essentialist Essence of Necessity.} is metaphysically reasonable. One reason to endorse (EEN), of course, is from the plausibility of essentialism. We established that essentialists should take (NP) to obtain essentially, and other candidates for the subject of the essentiality fails. Thus, the only remaining way to insist on essentialism is to endorse (EEN); and since essentialism implies the thesis, the thesis would be at least as reasonable as essentialism. I will not explore the theoretical motivations for essentialism here, for I do not aim in this paper to argue for the essentialist answer against other alternatives. Yet those motivations would also motivate (EEN). 
\par
There is another way to approach (EEN) with \citeA{Lowe2008-LOWTNO}'s slogan that \textit{essence precedes existence}. The idea, as I conceive, is as follows. To ask (with ontological significance, or in so-called ``ontologese''\footnote{\citeA{Sider2009-SIDOR-2}; also see \citeA{Dorr2005-DORWWD}}) whether $X$s exist, or even to deny the existence of $X$s we must at least partially know what $X$s are, that is, what the nature of $X$ is (also see \citeauthor{Fine2017-FINNM-2}, 2017). Thus, if the denial of the existence of $X$s is sensible, we could ascribe some facts to be essential of $X$s - or at least, we cannot deny that $p$ is essential of $X$s merely by the nonexistence of $X$s.
\par
It is common to think that there are varieties of necessity, regardless of whether there is a single fundamental kind of necessity where any other necessities are derivative. And all various kinds of necessities will have different essence, and make different facts obtain essentially. For example, it is plausible that given the nomic necessity operator $\Box^N$ and the predicate ``$NL$'' for ``\textit{to be a fact of natural laws}'', $E_{\Box^N} \forall p (NL(p) \supset \Box^N p)$. With the Lowean slogan, we do not have to worry about whether the nomic necessity is fundamental, or even \textit{natural} in the ontological sense. Now suppose we carve out - not necessarily along the joints of nature - a new kind of necessity: the \textit{essentialist necessity}. Let ``$\Box^E$ be the essentialist necessity operator. Now suppose that the partial nature of \textit{to be essentialistically necessary} is captured by (EEN)\textsuperscript{E}: $E_{\Box^E} \forall p(E_F p \supset \Box^E p)$\footnote{We can even suppose that $E_F p \bm{\leftrightarrow}_p \Box^E p$, viz., \textit{to obtain essentially} is the constitutive necessary and sufficient condition for \textit{to be essentialistically necessary}.}. If the realized essence fact of \textit{to be nomically necessary} is plausible, so should (EEN)\textsuperscript{E}. And from (EEN)\textsuperscript{E} to get (EEN) one only need to argue that the metaphysical necessity (``necessity'' for short throughout the paper) just is what we previously called essentialist necessity.\footnote{The general idea of this scheme is drawn from reading \citeauthor{Fine2020-Reply} (2020, p.463-4), although I do not have confidence whether this aligns with Fine's intention. At least \citeA{Romero2023-CAROEN} seems to interpret it in the opposite way, where Fine is taken to be committed to something similar to (MEE).} 
\par
Again, to argue for the equivalence between two kinds of necessity would require another twelve or more pages, and definitely exceed the scope of this paper. Still, the tenability of (EEN)\textsuperscript{E} would at least indirectly show that (EEN) would not be metaphysically unlikely like (MEF). If one can acknowledge (EEN)\textsuperscript{E} but not (EEN), one should provide reason why metaphysical necessity is crucially different from the essentialist necessity;\footnote{Notice that although I argue that (EEN) is forced upon essentialists, the converse does not hold. For example, the thesis that both necessity and essentiality is fundamental is compatible with (EEN).} and there seem no quick argument for it. At least on the surface-level, there are abundant similarities between them. For example, both take facts to be the subject, both are alethic necessities, both are non-epistemic, and so on. 
\par
And I further argue that essentialists can answer the challenge of gap by appealing to (EEN). To briefly reiterate, their main criticism was that once we stick to the notion of essence that is \textit{not} understood in terms of necessity, there seems to be a gap between essence and necessity. I agree with Romero (2019) that essentialists cannot take facts obtaining essentially to be \textit{ab initio} necessary. Nevertheless, there are many philosophers with  the essentialist intuition that it is in some sense ``obvious'' that facts obtaining essentially are necessary. There certainly is some kind of frustration when essentialists are confronted with questions like “why are facts of objectual essential necessary?” or “why an object necessarily has the property that defines what the object is?” This is, I suggest, same kind of frustration when one is confronted with questions like ``why is water $\text{H}_2 \text{O}$?'' or ``why does Socrates' singleton set have Socrates as member?''. As \citeA{Glazier2017-GLAEE-2} argues, we can simply answer: ``that is just what water is'' or ``that is just what the set is''. Turn to the question``Why are facts obtaining essentially necessary?'', I suggest answer essentialists can give is: ``that is just what to be necessary is!''. And to give such an answer is to accept (EEN). 
\par
Thus, although essentialists must endorse (EEN) with independent reason, this simultaneously indicates the possible reply to the challenge of gap. Essentialist can reply that it is misleading to ask whether \textit{to obtain essentially} has a suitable nature to be the source of necessity. Rather, it is the nature of \textit{to be necessary} is intimately related to \textit{to obtain essentially}, and that provides the metaphysical basis for (NP) to hold. Mackie is correct to point out that essentialists cannot deliver a modal rabbit (\textit{necessity} of facts obtaining essentially) out of a non-modal hat (non-modal conception of essence). Yet essentialists simply do not have to. What they have to, and are able to do is to deliver an essentialist rabbit (necessity of \textit{facts obtaining essentially}) out of an essentialist hat (essentialist conception of metaphysical necessity). 
\par
Furthermore, some relation between essence and necessity assumed by essentialist to vindicate (NP) could be inferred from (EEN). For example, let us revisit Hale's assumption that given $E_\alpha \Phi \alpha$, $\Diamond \lnot \Phi \beta \supset \alpha \neq \beta$. Suppose for contradiction that $\alpha = \beta$ and therefore by substitutability $\lnot \Box \Phi \alpha$. Now, from (EEN) we infer $\forall p (E_F p \supset \Box p)$, and thus $E_\alpha \Phi \alpha \supset \Box \Phi \alpha$. Finally, by the assumption $E_\alpha \Phi \alpha$ and modus ponens we get $\Box \Phi \alpha$. Contradiction. Therefore, there can be no $\beta$ such that $\alpha = \beta \land \lnot \Box \Phi \beta$, which ensures that if $\Phi \alpha$ obtains essentially of $\alpha$, then $\Phi$ is a metaphysically necessary condition for something to be identical to $\alpha$. 
\par
And the necessity of generalized identity. Although (EEN) is sufficient to secure the necessity of generalized identity\footnote{If we agree that $p \equiv q$ indicates that $q$ is essential of $p$, since the fact that $q$ is essential of $p$ would obtain essentially of $p$, we could infer that $E_p p \equiv q$. And by (EEN) \textit{to be necessary} is such that whatever obtains essentially is necessary, $p \equiv q$ is necessary.}, if one is sympathetic to \citeA{Correia2017-CORGEA} in the same spirit with (EEN) one can instead adopt $E_\Box ((\phi \equiv \psi) \supset \Box(\phi \equiv \psi))$. That is, \textit{to be necessary} is essentially such that if a generalized identity actually holds, it holds necessarily. It is again a mistake to assume if the source of the necessity of generalized identity must be found from the essence of \textit{to be generically identical to}, or what \textit{to be generically identical to} generically identical to. Instead, essentialists can simply answer all generalized identities are necessary because that is just \textit{to be necessary} is. 


%-----------------------------------------
\section{Circularity}
%-----------------------------------------
I argued in the previous section that essentialists should accept (EEN). This was done by showing that any other alternative is incompatible with essentialism. But to say that it is the only viable option for essentialism is not to say that it is tenable. After all, essentialism might also be incompatible with (EEN); and together with the argument I have presented, it would complete a trilemma that undermines essentialism.
\par
It seems clear that (EEN) is not \textit{logically} inconsistent with essentialism. To say what it is to be necessary is such and such is one thing, and to say the necessity is explained by such and such factor is another. And I also argued in \S 3.3 that (EEN) by itself is not metaphysically awkward, which it was in the case of (MEF). But it could be the case that the thesis seriously undermines the theoretical motivation of essentialism. The reason for rejecting (NP) obtaining essentially of \textit{to obtain essentially} was not logical inconsistency. The argument against it was that it gives a verdict on priority which is inconsistent with essentialism. 
\par
As I argued in \S 3.3., (EEN) would not face the same problem. But perhaps there are other challenges. Romero (2019) states six constraints on the notion of explanation which the `explanation' of modality by essence must satisfy: 

\begin{tabular}{p{16cm}}
\textbf{Why} The explanation must answer to a `why' question. \\
\textbf{Acyclicity} The explanation should not be circular. \\
\textbf{Fundamentality} The explanans should be metaphysically more fundamental than the explanandum. \\
\textbf{Ontic and Factive} It must be the case that the truth of the explanans, not our concept or theory of it, explains the truth of the explanandum. \\
\textbf{Structural} The explanandum need not be an object or a property. \\
\textbf{Objectivity} The explanation should be objectively valid. (p.124) \\
\end{tabular}
\\

It is already shown that (EEN) does not make essentialism infringe \textbf{Fundamentality}. Essentialism with or without (EEN) also easily satisfies \textbf{Why}, \textbf{Ontic and Factive}. It answers the question ``\textit{Why} is $p$ necessary?'', by appealing to the \textit{facts} about essence of things, not our concept of essence. Essentialists would certainly hold that if the essentialist account on necessity is valid, then it is \textbf{objectively} valid. Finally, \textbf{Structural} is not a negative constraint that one can infringe. 
\par
However, whether essentialism with (EEN) is compatible with \textbf{Acyclicity} is perhaps not so obvious. Some might feel uncomfortable with the idea that the nature of \textit{to be necessary} validates the explanation from \textit{to obtain essentially} to \textit{to be necessary}. But if it is indeed circular, which exact part of the picture is viciously circular?
\par
First, notice that I do not try to explain the necessity of $p$ by (EEN). The explanatory work is done by the fact that $p$ obtains essentially, and (EEN) is not involved as the source of the necessity of $p$. However, it is even unclear why including (EEN) as the source would be viciously circular. Suppose I am wrong, and every source of necessity must include (EEN). $p$ is necessary because $p$ obtains essentially \textit{and every such fact is necessary in virtue of the nature of to be necessary}. What is wrong with this kind of explanation? My point can perhaps be made by a simple analogy. If someone asks me why 3 is a prime number, I can simply answer ``because 3 is a natural number that is not a product of two natural numbers smaller than itself.'' However, if someone thinks this is incomplete, then I can add: ``... \textit{and} every such natural numbers is prime in virtue of the nature of \textit{to be prime}''. It would be obvious that neither explanation for the fact that 3 is prime is viciously circular. 
\par
Perhaps another kind of charge can be made by pinpointing the particular necessity of (NP). As I argued in \S 3, essentialists must admit that (NP) is necessary, and the fact that $\Box$(NP) would be grounded by (EEN), that is, $E_{\Box}$(NP). And this would seem somewhat problematic. If in explaining why (NP) is necessary essentialists must appeal to the nature of \textit{to be necessary}, this would seem circular.
\par
However, I still do not see anything problematic here. To generalize, someone who finds the situation problematic would assume that if $\phi$ consists of $X$, then $E_X \phi$ cannot ground $X(\phi)$. But this principle is not convincing. For example, suppose a quasi-God (`$Q$' below) knows all and only facts obtaining essentially, and its knowing a fact obtaining essentially is metaphysically grounded by that fact obtaining essentially. And suppose it is essential of \textit{to know} that if someone knows that $p$, it is the case that $p$, that is, $E_K \forall p (\exists s K(s, p) \rightarrow p)$. It would then follow that $E_K \forall p(K(Q, p) \rightarrow p)$ in the consequential sense. Now, by the assumption on $Q$ it would be the case that $K(Q, \forall p(K(Q, p) \rightarrow p))$, and it is partially grounded by $E_K \forall p(K(Q, p) \rightarrow p)$. It seems clear there is nothing viciously circular here. But this is exactly the instance being discussed: namely, there being a $\phi$ containing $X$ such that $E_X \phi$, where $E_X \phi$ grounds $X(\phi)$. 
 
%-------------
\section{Conclusion}
%-----------------
I argued that essentialists must accept (EEN), and by accepting (EEN) they could respond to the challenge of gap. The challenge points out essentialists must adopt (NP), but with the non-modal conception of essence one cannot infer (NP) from the notion of essence. Independently, I further argue that essentialists must also adopt $\Box$(NP), and the only way to account for $\Box$(NP) is to accept (EEN). And it turns out (EEN) can also be utilized to reply to the challenge of gap. Indeed (NP) is not explicable from the non-modal conception of essence; but it is explicable from the essentialist conception of metaphysical modality.
\par
What I did \textit{not} argue in this paper is that essentialism is superior to its competitors. My aim in this paper has been to show that the particular challenge to essentialists can be met coherently via what essentialists must endorse anyway. The aim of the paper is defensive: although what the challenge from gap points out is correct, it is not strong enough to defeat essentialism. It rather helps essentialists to clarify what comes together with essentialism. \textit{If} essentialism is correct, then so is (EEN). Any argument in favor of essentialism is also an argument for (EEN). Furthermore, if one agrees that \textit{essence precedes existence} \cite{Lowe2008-LOWTNO}, it illuminates another argumentative strategy for essentialists. But of course, the contrapositive of the conclusion also holds. If (EEN) is false, so is essentialism; thus any criticism against (EEN) would thereby be a criticism of essentialism. However, fleshing out new arguments for essentialism, or dealing with possible criticisms of (EEN) would exceed the scope of this paper. 
% --- BIBLIOGRAPHY --
\bibliography{refs} 
\nocite{Wildman2018-WILATR-7}
\nocite{Fine2002-FINVON}
\nocite{Bovey2022-BOVOTN}
\end{document} 